\section{引戻し量子化の非完全交叉変形環}
\label{sec:pullback-fixed-hulls-v37}

本節では、二つの独立な障害変数を持つ固定量子圏を構成する。
接線共鳴族の一変数切断環と異なり、一般の引戻し量子化には、
二次元の被約変形空間と、Cohen--Macaulay型二の非被約厚みが同時に残る。
入力は、公表稿IX\cite{VKPoisson34}の共点三直線における全局所環と、
その引戻し族の微分Ore量子化である。

\subsection{対象と主定理}

$m\ge3$、$N_m=\binom{m+2}{2}$ とする。$x,y,z$ の重みを一、$w$ の重みを $m$ とし、
\[
 c=xy(x+y),\qquad L=Ax+By,\qquad f=u x^{m+1}
\]
を取る。$\mathbf C[x,y,w]$ 上の重み一次の導分を
\begin{align*}
 \delta(x)&=\frac{c_y}{3}+\frac{m}{3(m+3)}xL,\\
 \delta(y)&=-\frac{c_x}{3}+\frac{m}{3(m+3)}yL,\\
 \delta(w)&=f-\frac{m}{m+3}Lw
\end{align*}
と定める。$\operatorname{div}\delta=0$ であり、
\[
 S_m(A,B,u;s)=\mathbf C[x,y,w][z;s\delta]
\]
はHilbert級数 $((1-t)^3(1-t^m))^{-1}$ のAS正則4-Calabi--Yau代数である。
この構成では $x,y,w$ は互いに可換で、$[z,v]=s\delta(v)$ である。
標準全列の有限代数を
\[
 \mathcal A_m(A,B,u;s)
 =\operatorname{End}\left(\bigoplus_{i=0}^{m+2}\mathcal O_{S_m}(i)\right)
\]
と書く。

\begin{theorem}[二変数の非完全交叉厚みを持つ固定量子圏]
\label{thm:pullback-fixed-hull-v37}
各 $m\ge3$ に対し、$(2,5)$ を含む非空Zariski開集合
$U_m\subset\mathbf A^2_{A,B}$ が存在する。
$(A,B)\in U_m$、$u,s\ne0$ とすると、$\mathcal A=\mathcal A_m(A,B,u;s)$ の
miniversal環は非標準的に
\[
 \boxed{R_{\mathcal A}\simeq
 \frac{\mathbf C[[t_1,t_2,a,b]]}{(a^{m+2},ab,b^2)}}
\]
となる。また
\[
 \dim HH^\bullet(\mathcal A)=
 \begin{cases}
 (1,2,4,9)&m=3,\\
 (1,2,4,6)&m\ge4,
 \end{cases}
 \qquad HH^{>3}(\mathcal A)=0.
\]
特に、固定された $A=2,B=5,u=s=1$ は、全ての $m\ge3$ に対して結論を満たす。
\end{theorem}

一つの明示的な十分条件として、次の開集合を $U_m$ に取れる。
\[
 \kappa_1=1+\frac m3(1-B),\quad
 \kappa_2=1+\frac m3(1+A),\quad
 \kappa_3=1+\frac m3(1-A+B),
\]
\begin{align*}
 \Delta_m(A,B)={}&(B-2A)(A-2B)(A+B)
                  \kappa_1\kappa_2\kappa_3\,[m(B+2)+6]\\
 &\hspace{12mm}\cdot\prod_{k=0}^{m}[m(B+2)-3k],\\
 U_m={}&\{\Delta_m\ne0\}\setminus
       \bigcup_{\substack{\nu\in\mathbf Z_{>0}^3\\|\nu|=m+3}}
       \{\kappa=\nu\}.
\end{align*}
これらは共点基底切片、全ての次数 $m$ の移動の単射性、
残す四行のSchur切片、および三角形側の非共鳴を保証する。
例えば最高重量で消す二行の小行列式は
\[
 -\frac{[m(1-B)+3][m(B+2)+6]}9
\]
である。他の重量では標的への全射性が $\kappa_i\ne0$ から従うため、
各点の近傍で可逆小行列を選べる。
$(2,5)$ では $\kappa_1<0$、$\kappa_2,\kappa_3>0$、
$m(B+2)-3k=7m-3k>0$ であり、全ての条件が成立する。
この $U_m$ は十分な範囲を明示するもので、全例外パラメータの分類ではない。

\subsection{非零混合点の全Poisson局所環}

\begin{lemma}[残余混合係数を反転した局所環]
\label{lem:pullback-nonzero-local-v37}
上の開集合を十分小さく取る。$u_0\ne0$ とし、
\[
 \pi=(c,D,[F]),\qquad
 D=(Ax+By)\partial_z,\qquad F=u_0x^{m+1}\partial_z
\]
を重み付きPoissonスキーム $\mathscr P_m$ の点とする。このとき
\[
 \widehat{\mathcal O}_{\mathscr P_m,\pi}
 \simeq
 \frac{\mathbf C[[z_1,\ldots,z_{N_m+9},a,b]]}
 {(a^{m+2},ab,b^2)}.
\]
\end{lemma}

\begin{proof}
公表稿IXの共点切片では、六つの $\mathrm{SL}_3$ 軌道変数を除いた基底は
\[
 \mathbf C[[A-A_0,B-B_0,a,b]]/(ab,b^2)
\]
であり、混合写像は $M_0+aM_1+bM_2$ となる。
同稿の全重みの核計算により、次数 $m$ の $w$ 移動は
混合核の階数 $N_m$ の自由直和因子をなし、残る一方向は
$x^{m+1}\partial_z$ で生成される。
Schur消去後の残余方程式は、単元倍を除いて正確に
\[
 a^{m+2}u=0
\]
である。$b$ の一次項は重量の差によって消え、他の $b$ 項は $ab=b^2=0$ により消える。

この消去を、$F=0$ の完備局所環へ非零の複素数 $u$ を代入する操作と解釈してはならない。
実際、基底 $(c,D)$ の座標変換は $F$ に依存せず、
混合写像のSchur消去と源の基底変換は全て $F$ の係数について線形である。
従って、基底方向の完備化と、$F$ の係数に関する多項式環との上で
同じ消去が成立する。任意の有限の $F$ の値で改めて完備化してよい。
$u=u_0+v$ と置けば $u_0+v$ は単元なので、残余イデアルは $(a^{m+2})$ となる。
六つの軌道変数、$A,B$、$N_m$ 個の移動変数、$v$ を合わせると、
滑らかな変数は $N_m+9$ 個である。
\end{proof}

\begin{lemma}[Poisson安定化群の二方向]
\label{lem:pullback-poisson-stabilizer-v37}
補題\ref{lem:pullback-nonzero-local-v37}の点における
重み付き座標変換群の軌道次元は $N_m+7$ であり、
安定化Lie代数は
\[
 \mathbf C\,x\partial_z\oplus\mathbf C\,y\partial_z
\]
である。
\end{lemma}

\begin{proof}
座標変換群のLie代数は、Euler場を除いて
$\mathfrak{sl}_3$、$w\partial_w$、$g_m\partial_w$ からなり、次元は $N_m+9$ である。
三つの相異なる共点直線の積 $c=xy(x+y)$ の $\mathrm{SL}_3$ 軌道は六次元である。
その安定化Lie代数は $x\partial_z,y\partial_z$ の二方向であり、
両者は $D=L\partial_z$ も保存する。

公表稿IXの核計算では、全ての次数 $m$ の $w$ 移動の像が $N_m$ 次元であり、
$F_{\mathrm{ex}}=x^{m+1}\partial_z$ の類はその像に入らない。
特に二変数移動では、$V=(H_c+LE_{xy})/3$ として
\[
 V(x^{m-k}y^k)=
 \left(\frac{m(A+1)}3-k\right)x^{m+1-k}y^k+
 \left(\frac{m(B+2)}3-k\right)x^{m-k}y^{k+1}.
\]
$(A,B)=(2,5)$ では後者の係数は $7m/3-k>0$ なので単射であり、
余核は $x^{m+1}$ で生成される。非二変数移動の単射性は重量ごとの核の表示から従う。
$u_0\ne0$ なので $w$ の倍率がさらに独立な一方向を与える。
従って全軌道の次元は少なくとも $6+N_m+1=N_m+7$ である。

一方、発散零代表 $\widehat\pi=\partial_z\wedge\delta$ は
$x\partial_z,y\partial_z$ によって保存される。
例えば $[x\partial_z,\delta]=-\delta(x)\partial_z$ であり、
これを $\partial_z$ と外積すれば零となる。
従って安定化群の次元は少なくとも二で、軌道次元の上界も $N_m+7$ となる。
上下界が一致する。
\end{proof}

\subsection{Artin厚みを保つ形式移送と軌道切片}

次の二補題は、一変数のArtin環に限らず使える。

\begin{lemma}[定数Artin厚みを持つ錐の弧]
\label{lem:artin-cone-arc-v37}
$k$ を標数零の体、$T$ を剰余体 $k$ の有限局所 $k$ 代数とする。
斉次アフィン錐 $P$ の点 $\pi$ において
\[
 \widehat{\mathcal O}_{P,\pi}\simeq T[[z_1,\ldots,z_d]]
\]
とする。$p(h)=h\pi+O(h^2)$ が $P(k[[h]])$ の弧なら、
$K=k((h))$ 上で
\[
 \widehat{\mathcal O}_{P_K,p_K}
 \simeq(T\otimes_kK)[[z_1,\ldots,z_d]].
\]
\end{lemma}

\begin{proof}
錐のスカラー作用で $p(h)$ を $p(h)/h=\pi+O(h)$ に移す。
形式弧に沿う移送補題\ref{lem:arc-translation-v35}を適用する。
$k[[h]]$ は被約なので、$T$ の極大イデアルは全て零へ写る。
従って弧は滑らかな $z$ 方向だけを動かし、その値を形式的に平行移動すればよい。
\end{proof}

\begin{lemma}[滑らかな軌道方向だけを除く切片]
\label{lem:artin-orbit-slice-v37}
有限局所 $k$ 代数 $T$ の剰余体を $k$ とし、
$R=T[[z_1,\ldots,z_d]]$ とし、$\operatorname{char}k=0$ を仮定する。
滑らかな代数群の作用に対し、閉点の軌道の次元を $q$ とする。
このとき、軌道に横断するminiversal形式切片の環は
\[
 T[[z_{q+1},\ldots,z_d]]
\]
と取れる。
\end{lemma}

\begin{proof}
群のLie代数から、軌道接写像が単射となる $q$ 方向 $v_1,\ldots,v_q$ を選ぶ。
標数零なので形式指数 $\exp(\sum t_iv_i)$ を使え、その逆は $\exp(-\sum t_iv_i)$ である。
軌道は被約スキームの像であるから、$T$ の冪零元は軌道上で零となる。
従って軌道の接空間は $z$ 方向に含まれる。
$z$ の線形座標変換により、最初の $q$ 個の $z$ の軌道接写像を可逆にできる。
切片を $z_1=\cdots=z_q=0$ と定める。

任意の形式的な点について、群の上の $q$ 座標を使って、作用後の最初の $q$ 個の
$z$ 座標を零にする方程式を解く。そのJacobianは閉点で可逆なので、
形式逆関数定理により一意な解を得る。
従って全形式近傍は、この切片と $q$ 個の軌道変数との積になる。
切片の接空間は軌道接空間を割った空間であり、miniversal性を持つ。
\end{proof}

体拡大から元の複素係数へ同型型を戻す際には、次の初等的な補題を使える。
これは有限次元だけの環に限らない。

\begin{lemma}[複素係数の形式同型の下降]
\label{lem:complex-formal-descent-v37}
$R,R'$ を剰余体 $\mathbf C$、有限埋込次元の完備Noether局所
$\mathbf C$ 代数とする。ある拡大体 $K/\mathbf C$ に対して、$K$ 代数として
\[
 R\widehat\otimes_{\mathbf C}K
 \simeq R'\widehat\otimes_{\mathbf C}K
\]
なら、$R\simeq R'$ である。同型の選択は非標準的である。
\end{lemma}

\begin{proof}
$R=\mathbf C[[X_1,\ldots,X_e]]/(f_1,\ldots,f_r)$、
$R'=\mathbf C[[Y_1,\ldots,Y_{e'}]]/(g_1,\ldots,g_{r'})$ と表示する。
全ての $f_i,g_j$ の係数は可算個なので、これらを含む可算部分体
$k_0\subset\mathbf C$ を取れる。
与えられた同型と逆同型について、各変数の像を形式冪級数に持ち上げる。
定義関係が定義イデアルへ入ることと、二つの合成が恒等写像になることを表す
有限個のイデアル所属の係数冪級数も選ぶ。
これら全ての係数を $k_0$ へ添加して得る体 $K_0\subset K$ は可算である。

$\mathbf C$ の $k_0$ 上の超越次数は非可算なので、$K_0/k_0$ の
超越基底を $\mathbf C$ の代数的独立な集合へ送り、代数的部分の埋込みを延長することで、
$k_0$ 上の体の埋込み $K_0\hookrightarrow\mathbf C$ が存在する。
これを、同型・逆同型および所属を表す全係数に適用する。
各係数恒等式は有限個の体演算からなるので保たれ、$f_i,g_j$ の係数は固定される。
変数の像を $\mathbf C$ 線形に延長すれば、元の二つの完備局所環の間の
互いに逆な同型を得る。
\end{proof}

\subsection{量子側の $HH^1$ と主定理の証明}

Ore族の積はPBW基底上で
\[
 a\star_s b=\sum_{r\ge0}\frac{s^r}{r!}
                  (\partial_z^r a)(\delta^r b)
\]
と書ける。各対について有限和なので、$s=h$ は有限代数の多項式的な量子化弧を与える。
$h\ne0$ 上では、$x,y,w$ を固定して $z\mapsto hz$ とすることで
$S_m(A,B,u;h)\simeq S_m(A,B,u;1)$ を得る。
また $w\mapsto uw$ により、非零 $u$ も一へ規格化できる。

\begin{lemma}[量子側の二つの外的導分]
\label{lem:pullback-quantum-HH1-v37}
主定理の仮定の下で $\dim HH^1(\mathcal A)=2$ である。
生成元は、座標環の自己同型
\[
 z\longmapsto z+tx,\qquad z\longmapsto z+ty
\]
を $t=0$ で微分した導分から得られる。
\end{lemma}

\begin{proof}
$x,y,w$ は可換なので、上の二つの置換はOre関係を保ち、逆も同じ形で与えられる。
次数も保つため、標準有限代数の全頂点を固定する自己同型を誘導する。
対応する導分を $D_x,D_y$ と書けば、生成元上で
\[
 D_x(z)=x,\quad D_y(z)=y,\qquad
 D_x(x)=D_x(y)=D_x(w)=D_y(x)=D_y(y)=D_y(w)=0.
\]
PBW多項式上の通常の $x\partial_z,y\partial_z$ と、そのまま同一視する必要はない。

冪等元を固定する内的導分は、対角のスカラー元から来る。
従って各次数一の隣接射空間にはスカラー倍で作用する。
上の二つの導分の非零線形結合は、$x,y$ を零へ送り $z$ を非零の一次式へ送るため、
内的ではない。この議論は $\mathbf C((h))$ 上でも成立する。
よって $HH^1$ の次元は少なくとも二である。

逆向きには、最初のHochschild差分補題\ref{lem:first-differential-v35}を使う。
可換基点では $HH^1=H^0(T_{X_m})$、一次差分は $[\pi,-]$ である。
補題\ref{lem:pullback-poisson-stabilizer-v37}からその核は二次元であり、
$HH^0=\mathbf C$ からの差分は零なので、generic fiberで $HH^1\le2$ となる。
非零係数間の同型によって全ての $s\ne0$ へ移せば結論を得る。
\end{proof}

\begin{proof}[定理\ref{thm:pullback-fixed-hull-v37}の証明]
$K=\mathbf C((h))$ とする。可換標準列代数のhullとPoisson錐の原点完備環の同型は、
HKRを接写像として持つ。従って、代数構造スキーム $\mathfrak A$ の可換点における
全完備局所環は、Poisson錐の原点環に $b_0$ 個の軌道変数を付け加えたものである。
ここで $b_0$ は可換点でのHochschild境界 $B_E^2$ の次元である。

Ore弧のPoisson座標は $h\pi+O(h^2)$ である。
補題\ref{lem:pullback-nonzero-local-v37}と
補題\ref{lem:artin-cone-arc-v37}、形式弧の移送から、generic fiberの全局所環は
\[
 \widehat{\mathcal O}_{\mathfrak A_K,\mathcal A_h\otimes K}
 \simeq
 \frac{K[[z_1,\ldots,z_{b_0+N_m+9},a,b]]}{(a^{m+2},ab,b^2)}
\]
となる。

中心はどの標準有限代数でもスカラーのみで、$\dim\operatorname{im}d^0=m+2$ は一定である。
可換点の $HH^1$ は $N_m+9$ 次元、量子側では補題により二次元なので、
基底変換軌道の次元は
\[
 b_\eta=b_0+N_m+7
\]
となる。補題\ref{lem:artin-orbit-slice-v37}を適用すれば、二つの滑らかな変数が残り、
generic fiberのhullは
\[
 \frac{K[[t_1,t_2,a,b]]}{(a^{m+2},ab,b^2)}
\]
である。

非零係数の同型は $\mathcal A_h\otimes K\simeq\mathcal A_1\otimes K$ を与える。
有限代数の形式切片は体拡大と可換なので、補題\ref{lem:complex-formal-descent-v37}により、
この同型型は $\mathbf C$ 上へ降下する。全ての非零 $s,u$ も同じである。

得られたhullの埋込次元は四なので $HH^2$ の次元は四である。
$HH^{>3}=0$ は、可換点での消滅、有限自由Hochschild複体の特殊化、
および非零係数間の同型による。
同じ有限複体のEuler特性は $m=3$ で $-6$、$m\ge4$ で $-3$ であり、
$HH^0=1,HH^1=2,HH^2=4$ を代入すれば表示した $HH^3$ の次元を得る。
\end{proof}

\begin{corollary}[二つの滑らかな方向の明示実現]
\label{cor:pullback-reduced-parameters-v37}
主定理の点 $(A_0,B_0,u,s)$ を固定する。
$A=A_0+p,B=B_0+q$ とするOre族は、hullの被約化
\[
 (R_{\mathcal A})_{\mathrm{red}}\simeq\mathbf C[[t_1,t_2]]
\]
を形式的にパラメータ付ける。すなわち、hullからこの族への写像を一つ選べば、
誘導される写像
$(R_{\mathcal A})_{\mathrm{red}}\longrightarrow\mathbf C[[p,q]]$
は同型である。
\end{corollary}

\begin{proof}
まずPoisson側の二つの接ベクトル $\partial_A\pi,\partial_B\pi$ は、
座標変換軌道の接空間を法として独立である。
実際、この線形結合を軌道接ベクトルと仮定すると、その $c$ 成分は零なので、
$\mathfrak{sl}_3$ 成分は $x\partial_z,y\partial_z$ の線形結合である。
両者は $D$ を固定し、$w$ 移動と倍率も $D$ を変えない。
一方、上の二つのパラメータは $D=(Ax+By)\partial_z$ を独立に変えるため、
線形結合の係数はともに零となる。

量子側への独立性の移送も確認する。
$O=\mathbf C[[h]]$ 上で可換点の可縮部分を除いたHochschild複体を取り、
次数一から二への微分を
\[
 d_h=h(D_1+hD_2+\cdots),\qquad D_1=[\pi,-]
\]
と書く。$HH^1$ の次元は、$\ker D_1$ とgeneric fiberの両方で二である。
次数零の微分は中心の次元が一定一であるため零なので、
$\operatorname{rank}_K(d_h/h)=\operatorname{rank}_{\mathbf C}D_1=N_m+7$ となる。
従って $d_h/h$ の像は $O$ 上の自由な直和因子である。
実際、$D_1$ の可逆小行列を消去すれば、残余行列は $K$ 上で階数零、
従って $O$ 上でも零となる。

Ore族の二つのKodaira--Spencerコサイクルは $h$ で割り切れ、
上の縮約複体における商の先頭項は $\partial_A\pi,\partial_B\pi$ である。
これらが $K$ 上のコホモロジーで従属なら、係数を一つの $h$ 冪で規格化し、
少なくとも一つを単元とした $O$ 係数の関係を取れる。
$d_h/h$ の像が直和因子なので、この関係を $h=0$ へ特殊化すると、
$\partial_A\pi,\partial_B\pi$ の非零線形結合が $\operatorname{im}D_1$ に入る。
これは前段と矛盾する。
非零 $h$ と一との同型は $A,B$ に依存しないので、固定された非零 $s$ でも
二つのKodaira--Spencer類は独立である。

形式族はhullの性質により $R_{\mathcal A}\to\mathbf C[[p,q]]$ を与える。
標的は被約なので、この写像はhullの被約化を通る。
主定理により被約化は二変数の形式冪級数環であり、
今証明した接写像の独立性から二変数間のJacobianが可逆である。
形式逆関数定理によって結論を得る。
\end{proof}

\subsection{変形環の構造と検証}

$n=m+2$ と置く。横断Artin環
\[
 T_m=\mathbf C[a,b]/(a^n,ab,b^2)
\]
の基底は $1,a,\ldots,a^{n-1},b$、socleは
$\langle a^{n-1},b\rangle$ である。従って長さは $m+3$、socle次元は二である。
主定理の環は $\mathbf C[[t_1,t_2]]$ 上の階数 $m+3$ の自由加群であり、
二次元Cohen--Macaulay環、型二、Hilbert--Samuel重複度 $m+3$ を持つ。
特にGorensteinではなく、完全交叉でもない。
関係イデアルのHilbert--Burch行列は
\[
 \begin{pmatrix}
 b&a^{n-1}\\
 -a&0\\
 0&-b
 \end{pmatrix}
\]
と取れる。高さ二に対して最小生成元は三個である。

二つの障害方向は異なる振舞いを示す。
$b$ 方向は $b^2=0$ によって二次で止まり、$a$ 方向の最初の純粋な障害は
$a^{m+2}=0$ によって次数 $m+2$ に現れる。
両者は二次関係 $ab=0$ でも結び付く。
一方、被約な形式基底への変形は $a=b=0$ を通り、二つの滑らかな方向を残す。

付属の有理数検算は、全標準有限代数の導分を矢印像から計算し、
Ore関係のLeibniz式の全係数を照合する。
$(A,B,u,s)=(2,5,1,1)$ における結果は次の通りである。
\[
 \begin{array}{c|r|r|r|r|r}
 m&\text{未知数}&\text{方程式}&\operatorname{rank}_{\mathbf Q}
 &\dim\operatorname{Der}_E&\dim HH^1\\\hline
 3&78&180&71&7&2\\
 4&102&234&94&8&2\\
 5&129&294&120&9&2\\
 6&159&360&149&10&2
 \end{array}
\]
別係数 $(3,7,2,-1/2)$ と可換基点も照合し、
二つの明示導分が内的導分から独立なことを行列上で確かめた。
$A,B$ の関係微分を二列追加した行列では、全ての非零係数例で
階数が正確に二増え、二つのKodaira--Spencer類の独立性も確認できる。
Poisson安定化群も別の多項式Lie微分行列で検査し、
$m=3,4,5,6$ の軌道階数 $17,22,28,35$ と二次元の核を得る。
これらの有限計算は、本文の全重みの構造的証明を補助する。

本節は、固定した標準有限代数、従って対応する射影導来圏の全結合変形を扱う。
無限次元の座標環の全非次数付き変形や、
パラメータ $(A,B)$ 間の射影圏の完全な同値分類を主張するものではない。
